% GENERATED FILE. Edit canonical lessons, then run npm run generate:books.
% canonical-source-hash: fnv1a64:874ad3cef4ecf3da
% canonical-lessons: KA-C13-dehada-bhagagalu, KA-S122-letter-ca, KA-S144-digit-five

\chapter{Head and Hand}
\label{ch:ka-head-and-hand}
\hlchaptermodality{13}

\begin{chapteropening}
\textbf{By the end of this chapter:} \emph{I can name the head and the hand in Kannada.}

Say \kn{ತಲೆ} and \kn{ಕೈ}, and hear how nearly unchanged they are across Tamil, Malayalam and Kannada.
\end{chapteropening}

\section[tale kai]{\kn{ತಲೆ}, \kn{ಕೈ} — head and hand, the shared Dravidian core}

\label{lesson:KA-C13-dehada-bhagagalu}



\begin{quote}
Three languages, nearly one word each — Kannada's basic body parts line up with Tamil and Malayalam almost perfectly.
\end{quote}

\subsection*{You'll want to know: The words}
\begin{itemize}
\raggedright
  \item \textbf{\kn{ತಲೆ}} (\emph{tale}) = \textquotedblleft{}\textbf{head}\textquotedblright{} — native Dravidian, matching Tamil \emph{talai} and Malayalam \emph{thala}.
  \item \textbf{\kn{ಕೈ}} (\emph{kai}) = \textquotedblleft{}\textbf{hand}\textquotedblright{} — native Dravidian, matching Tamil and Malayalam's \emph{kai} almost exactly.
\end{itemize}

\subsection*{Your turn}
\begin{itemize}
\raggedright
  \item \emph{Say it:} \textquotedblleft{}tale\textquotedblright{} — head
  \item \emph{Say it:} \textquotedblleft{}kai\textquotedblright{} — hand
  \item \emph{Say it:} the match — nearly identical across Tamil, Malayalam, Kannada
  \item \emph{Recall:} say \emph{sōmavāra maṅgaḷavāra budhavāra guruvāra śukravāra śanivāra bhānuvāra}
\end{itemize}

\begin{tcolorbox}[breakable,colback=teal!4,colframe=teal!35!black,title={Before you move on}]
What are Kannada's words for head and hand? (\textbf{\kn{ತಲೆ} tale, \kn{ಕೈ} kai}.) How closely do they match Tamil and Malayalam? (\textbf{Very closely} — all three share essentially the same native Dravidian words.)
\end{tcolorbox}
\section[ca]{\kn{ಚ} — one character, met inside words you already say}

\label{lesson:KA-S122-letter-ca}



\begin{quote}
Before the new one: \kn{ಸ} — what does it do?

One character this time. Just one — and you have been saying it for pages without knowing which mark on the page it was.
\end{quote}

\subsection*{Script you'll notice: \kn{ಚ}}
\textbf{\kn{ಚ}} — \emph{ca}.

It is a \textbf{consonant}, and in this script a consonant is never bare: it comes with an \emph{a} already in it. So it is not \emph{c}, it is \textbf{ca}.

You already say these, and every one of them has \kn{ಚ} somewhere inside it:

\begin{itemize}
\raggedright
  \item \textbf{\kn{ಚೆನ್ನಾಗಿ}} \emph{cennāgi} — well, nicely — and the reply \textquotedblleft{}\kn{ನಾನು} \kn{ಚೆನ್ನಾಗಿದ್ದೇನೆ}\textquotedblright{}
  \item \textbf{\kn{ಚೈತ್ರ} \kn{ವೈಶಾಖ} \kn{ಜ್ಯೇಷ್ಠ} \kn{ಆಷಾಢ} \kn{ಶ್ರಾವಣ} \kn{ಭಾದ್ರಪದ} \kn{ಆಶ್ವಯುಜ} \kn{ಕಾರ್ತೀಕ} \kn{ಮಾರ್ಗಶಿರ} \kn{ಪುಷ್ಯ} \kn{ಮಾಘ} \kn{ಫಾಲ್ಗುಣ}} — the twelve lunisolar months, which the months lesson will name
\end{itemize}

\subsection*{Writing: \kn{ಚ} — copy what you see}
Put your pen on \kn{ಚ} and follow its line. Copy the shape you can see — slowly, and larger than it is printed.

\begin{quote}
This book does not yet tell you \textbf{where to start the character or which way to travel}. That is a real thing, taught with real variation from school to school, and it is not written down here until it can be written down with a source. Copying what is in front of you needs no such source, and it is how the shape gets into your hand in the meantime.
\end{quote}

\subsection*{Your turn}
\begin{itemize}
\raggedright
  \item \emph{Look:} at these words, and find \kn{ಚ} in the ones that have it
\end{itemize}

\begin{quote}
\kn{ಚೆನ್ನಾಗಿ}  ·  \kn{ನಮಸ್ಕಾರ}
\end{quote}

\begin{itemize}
\raggedright
  \item \emph{Trace it:} \kn{ಚ} three times, saying \emph{ca} as you finish each one
  \item \emph{Look:} back at any page of this chapter and find \kn{ಚ} once more
\end{itemize}

\begin{tcolorbox}[breakable,colback=teal!4,colframe=teal!35!black,title={Before you move on}]
Which character is this — \kn{ಚ}? What sound does it carry? (\emph{\textbf{ca}}.) Name one word you already say that contains it.
\end{tcolorbox}
\section[5]{\kn{೫} — the digit for aidu}

\label{lesson:KA-S144-digit-five}



\begin{quote}
Before the new shape: \textbf{\kn{೪}} — which number is it?

One digit this time, and you already know its name.
\end{quote}

\subsection*{Script you'll notice: \kn{೫}}
\textbf{\kn{೫}} — \textbf{5}, the digit for \textbf{\kn{ಐದು}} (\emph{aidu}), \textquotedblleft{}five\textquotedblright{}.

Halfway through the digits. Notice how little there is to learn here compared with the letters: no vowel signs, no stacking, no virama.

\subsection*{Writing: \kn{೫} — copy what you see}
Put your pen on \kn{೫} and follow its line. Copy the shape in front of you — slowly, and larger than it is printed.

\begin{quote}
This book does not yet tell you \textbf{where to start the character or which way to travel}. That is taught with real variation from school to school, and it is not written down here until it can be written down with a source. Copying what you can see needs no such source.
\end{quote}

\subsection*{Your turn}
\begin{itemize}
\raggedright
  \item \emph{Look:} at the digits you have met, and name each one aloud
\end{itemize}

\begin{quote}
\kn{೧}  ·  \kn{೨}  ·  \kn{೩}  ·  \kn{೪}  ·  \kn{೫}
\end{quote}

\begin{itemize}
\raggedright
  \item \emph{Say it:} \textquotedblleft{}aidu\textquotedblright{} as you point at \kn{೫}
\end{itemize}

\begin{tcolorbox}[breakable,colback=teal!4,colframe=teal!35!black,title={Before you move on}]
Which number is \textbf{\kn{೫}}? (\textbf{5} — \emph{aidu}, \kn{ಐದು}.) Name every digit you have met so far, in order.
\end{tcolorbox}
